\section{六、电池、充电、PMIC 与 USB-PD 电源路径}

台式机从电源供应器获得相对稳定的 12 V、5 V 和 3.3 V；移动设备则必须在外部适配器、电池和系统负载之间连续切换。电池电压随荷电状态、温度和瞬时电流变化，外部接口还要协商允许的电压与功率。移动平台的供电问题因而不只是稳压，而是能源来源选择、充电限制、多路电源时序、剩余容量估计和掉电恢复的联合状态机。

\begin{pdfdiagram}[x=1cm,y=1cm]
  \node[hw block, minimum width=2.4cm] (usb) at (1.4,4.0) {USB-C 端口\\VBUS / CC};
  \node[hw control, minimum width=2.7cm] (pd) at (4.6,4.0) {USB-PD 控制器\\能力协商与保护};
  \draw[hw bus] (usb) -- (pd);

  \node[hw memory, minimum width=2.4cm] (bat) at (1.4,1.2) {电池组\\电芯 + 保护 FET};
  \node[hw control, minimum width=2.7cm] (chg) at (4.6,1.2) {充电器/电量计\\CC/CV 与状态估计};
  \draw[hw bus] (bat) -- (chg);

  \node[hw compute, minimum width=3.0cm, minimum height=1.8cm] (path) at (8.1,2.6) {电源路径管理\\限流、理想二极管\\系统轨与电池轨切换};
  \draw[hw bus] (pd.east) -- (6.2,4.0) -- (6.2,3.2) -- (path.west);
  \draw[hw bus] (chg.east) -- (6.2,1.2) -- (6.2,2.0) -- (path.west);

  \node[hw control, minimum width=2.7cm] (pmic) at (11.5,2.6) {PMIC\\buck/LDO/负载开关};
  \draw[hw bus] (path) -- (pmic);
  \node[hw state, minimum width=2.4cm] (soc) at (11.5,0.2) {SoC/DRAM/外设\\多电源域};
  \draw[hw bus] (pmic) -- (soc);
  \draw[hw return] (soc.west) -- (9.7,0.2) -- (9.7,1.45) -- (path.south);
  \node[hw label,anchor=east] at (9.55,0.85) {负载与温度反馈};
\end{pdfdiagram}
\diagramnote{移动设备的完整供电通路。USB-PD 控制器决定外部输入上限，充电器决定电池允许的电流和电压，电源路径管理器在适配器、电池与系统负载之间分配能量，PMIC 再生成各电源域。任何一级限流或失效都会改变后级可用功率。}

\topic{电池不是理想电压源}

锂离子电芯的开路电压只能粗略反映荷电状态；带载时，内阻产生 $I R$ 压降，低温和老化又会提高内阻。同一块电池在轻载时可能显示 20\% 剩余容量，CPU 突然进入高功耗状态后端电压却跌破关机阈值。平台需要同时测量电压、电流、温度和累计电荷，并用电芯模型估计\term{荷电状态}{State of Charge, SoC}与\term{健康状态}{State of Health, SoH}。

电池保护电路独立于操作系统工作。过充、过放、过流、短路或温度越界时，保护控制器关闭串联 FET，把电芯与系统断开。软件显示的“0\%”通常仍高于保护芯片的硬下限，以便在正常关机后保留恢复与时钟所需能量；真正触发硬保护后，设备可能必须接入合格充电源才能重新打开电源路径。

\topic{充电器按恒流、恒压和温度窗口推进}

锂离子充电通常先以恒流阶段提高电芯电压，再在达到目标电压后进入恒压阶段，让电流逐渐下降。充电器根据电池温度、适配器能力、连接器温度和系统负载动态降低电流。若系统正在消耗 15 W，而外部适配器只承诺 10 W，多出的 5 W 仍来自电池；界面显示“已连接电源”不等于电池一定在充电。

充电完成的可见边界是充电电流低于终止阈值并满足稳定时间，不是 VBUS 出现或电芯达到一次目标电压。拔出适配器后，电源路径管理器要在输出电容允许的时间内切到电池，避免系统轨跌破 SoC 的 brownout 门限。

\topic{USB-PD 先建立通信契约，再提高 VBUS 功率}

USB Type-C 的 CC 引脚先确定连接方向、供电角色和默认电流。USB-PD 控制器随后在 CC 上交换 Source Capabilities、Request、Accept 和 PS\_RDY 等消息。供电端只有在接受请求并完成电压切换后，才能宣布新的电源契约；受电端在 PS\_RDY 前不能按高电压预算启动负载。

协商失败时，端口应回到安全的默认电压或断开，而不是让电源级停在未知状态。过压、过流、连接器过温与线缆能力不足分别由电源开关、采样电阻、温度传感器和电子标记线缆信息约束。USB 数据链路是否枚举成功与 PD 功率契约是否成立是两条独立状态机，调试时必须分别读取。

\topic{PMIC 把一个能源入口展开成有依赖的电源树}

PMIC 集成多路 buck、LDO、负载开关、时钟和复位控制。处理器核心、I/O、DRAM、模拟和射频电路可能要求不同电压，并规定上电先后、斜率和掉电顺序。电源域 A 的复位只有在对应轨达到电压窗口、参考时钟稳定且依赖域就绪后才能释放；关闭时则先停止新事务、保存必要状态，再关时钟和电源。

\begin{longtable}{L{0.18\linewidth}L{0.32\linewidth}L{0.38\linewidth}}
\toprule
观察点 & 能证明的事实 & 不能据此推出的事实 \\
\midrule
VBUS 存在 & 端口上检测到外部电压 & PD 已协商高功率，或电池正在充电 \\\tablerowrule
PD PS\_RDY & 供电端宣布目标电压已建立 & 系统每一路 PMIC 输出均已稳定 \\\tablerowrule
电池电压正常 & 当前端电压高于保护下限 & 高负载和低温时不会发生压降关机 \\\tablerowrule
PMIC power-good & 对应电源轨进入窗口 & 依赖时钟、训练和固件初始化均已完成 \\\tablerowrule
操作系统关机完成 & 软件已停止普通工作 & 掉电保持电容已放完，所有外设均无电 \\
\bottomrule
\end{longtable}

\topic{brownout 恢复必须从最早失效电源域开始}

电压短暂跌落可能只复位一部分器件：CPU 已复位，外设仍保留旧队列；PMIC 状态机已重启，DRAM 自刷新仍保存旧内容；存储控制器缓存尚未写介质。平台应使用独立 brownout 检测器锁存失效原因，并以全局代际号使跌落前的中断、DMA 完成和队列项失效。恢复时先确认能源来源和电源轨，再恢复时钟、复位、内存与外设，不能从“软件最后执行到哪里”直接继续。
